\documentclass[11pt]{article}

\usepackage[a4paper,margin=2.5cm]{geometry}
\usepackage{amsmath,amssymb,amsthm}
\usepackage{mathrsfs}
\usepackage[colorlinks=true,linkcolor=blue,citecolor=blue,urlcolor=blue]{hyperref}
\usepackage{booktabs}
\usepackage{enumitem}
\usepackage{natbib}
\usepackage{multirow}

\theoremstyle{plain}
\newtheorem{theorem}{Theorem}[section]
\newtheorem{conditional}[theorem]{Conditional Proposition}
\newtheorem{proposition}[theorem]{Proposition}
\newtheorem{corollary}[theorem]{Corollary}
\newtheorem{lemma}[theorem]{Lemma}
\theoremstyle{definition}
\newtheorem{definition}[theorem]{Definition}
\newtheorem{axiom}[theorem]{Axiom}
\newtheorem{hypothesis}[theorem]{Hypothesis}
\newtheorem{interpretation}[theorem]{Interpretation}
\theoremstyle{remark}
\newtheorem{remark}[theorem]{Remark}

\newcommand{\Zphi}{\mathbb{Z}[\varphi]}
\newcommand{\Fix}{\mathrm{Fix}}
\newcommand{\Ical}{\mathcal{I}}
\newcommand{\Ocal}{\mathcal{O}}
\newcommand{\Ccal}{\mathcal{C}}
\newcommand{\Cph}{C_\varphi}
\newcommand{\Xcal}{\mathcal{X}}
\newcommand{\catC}{\mathcal{C}}

\title{Cosmic Self-Pruning and Frame Recurrence in Closure Geometry\\[6pt]
\large A Structural Extension of the Bounded Reference Frame Framework}
\author{%
  Lee Smart\\[4pt]
  \textit{Institute of Vibrational Field Dynamics}\\[2pt]
  \texttt{contact@vibrationalfielddynamics.org}\\[2pt]
  \texttt{@vfd\_org}%
}
\date{May 2026}

\begin{document}

\maketitle

\noindent\textbf{Status:} Preprint (not peer-reviewed). Structural extension paper. Companion to \citet{SmartExistenceClosure} (core bridge paper) and \citet{SmartProcessingToPOV} (consciousness application). The core paper establishes existence-as-closure, bounded reference frames, and the candidate transition rule. The present paper develops the structural extension: frame dynamics (creation, destruction, merging, fission), frame recurrence (re-instantiation of frame attractors under compatible boundary conditions), cosmic self-pruning (closure-phase transitions across cosmological epochs), and the unified hierarchy mapping these life-cycle events onto the cascade rung structure. All structural propositions of the core paper are imported by reference. The propositions in this paper that depend on additional hypotheses (H-RP-1, H-RP-2, H-RP-3, H-rec, H-evolve, P-A) are labelled as \emph{Conditional Propositions} throughout; theorem-grade statements concern only the structural mathematics that follows from the closure-operator and symmetry axioms.

\begin{abstract}
We extend the closure-geometric framework of \citet{SmartExistenceClosure} to cover bounded reference frame life-cycle dynamics and the long-term evolution of the universal closure structure. Frame creation, destruction, merging, and fission are formalised as theorems on the closure dynamic $F_{t+1} = (I - \lambda \Cph) F_t + \kappa P_{\Sigma_\Ical}(B_t)$. A frame attractor is defined as a closure-stable $\sigma$-stable subspace $\mathcal{A} \subseteq C$ in the universal closure geometry, and frame compatibility $\kappa_{ab} = \dim(!_{\Ical_a}(\Ical_a) \cap !_{\Ical_b}(\Ical_b))$ is shown to be independent of local-time ordering. Under the conditional hypothesis H-rec (recurrence admissibility), a frame attractor may be locally instantiated by multiple bounded reference frames at distinct anchors; this is the framework's structural reframing of the phenomenology classically called \emph{reincarnation}, given here in the precise vocabulary \emph{frame recurrence}. At the cosmological scale, we introduce a pruning operator $P = P_{\Fix(\tau)} \circ P_{\Fix(\Cph)}$ and a closure-phase transition map $U_{n+1} = C_{n+1}(P(U_n))$; under hypothesis H-evolve (time-dependent closure operator), cosmological evolution is shown to be non-periodic, distinguishing the framework's prediction from both heat-death termination and classical cyclic-universe models. The unified hierarchy theorem states that the same structural template (creation, persistence, destruction, possible recurrence, closure-phase transition) applies at every cascade rung, with rung-specific time scales. The universe at the $E_8$ totality rung is identified as a bounded reference frame in its own right, with its own life-cycle, structurally distinct from biological life-cycles. The propositions throughout are conditional on H-rec and H-evolve; the core paper's structural results hold independently.
\end{abstract}

%==================================================================
\section{Introduction and the imported framework}\label{sec:intro}
%==================================================================

\subsection{Position relative to the core paper}

The companion paper \citet{SmartExistenceClosure} establishes the structural framework:

\begin{itemize}[leftmargin=2em]
\item A closure operator $\Cph = L_{V_{600}} + \varphi^{-2} I$ on the 600-cell substrate $V_{600}$.
\item A Galois reflection $\tau$ on $\Zphi$-icosian coordinates with $\dim \Fix(\tau) = 94$ unconditional.
\item A bounded reference frame $\Ical$ defined as a finite local projection of the substrate equipped with a self-referential constraint, admitting a per-frame zero-line $\Sigma_\Ical \subseteq \Fix(\tau)$.
\item The local-into-universal embedding $!_\Ical(v) = \pi_\Ical(P_{\Sigma_\Ical}(v))$ from each bounded reference frame to the universal closure structure $C = U \cap \Fix(\tau)$.
\item The candidate transition rule: a bounded local system supports point of view when its internal state closes under recursive update, producing non-trivial $\Sigma_\Ical$, conditional on the Access Principle (P-A).
\end{itemize}

The present paper develops the dynamical extension: frame life-cycle events, recurrence, and cosmic-scale phase transitions. All structural results of this paper rest on the closure operator's spectral properties and the symmetry $\tau$'s involutive structure; new biological identifications are explicitly conditional.

\subsection{Scope of this extension paper}

This paper does:

\begin{itemize}[leftmargin=2em]
\item Formalise frame creation, destruction, merging, fission as structural theorems on the closure dynamic (\S\ref{sec:frame-dynamics}).
\item Define frame attractors and frame compatibility, and prove time-independence of compatibility (\S\ref{sec:frame-recurrence}).
\item State the conditional proposition for frame recurrence under H-rec (\S\ref{sec:frame-recurrence}).
\item Define the pruning operator and closure-phase transition map; prove non-periodic recurrence (\S\ref{sec:cosmic-pruning}).
\item State the conditional proposition for closure-cascade cosmology under H-evolve (\S\ref{sec:cosmic-pruning}).
\item State the unified hierarchy theorem: structural life-cycle universality across cascade rungs (\S\ref{sec:hierarchy}).
\item Identify the universe at the $E_8$ totality rung as a bounded reference frame with its own life-cycle (\S\ref{sec:hierarchy}).
\end{itemize}

This paper does \textbf{not}:

\begin{itemize}[leftmargin=2em]
\item Reproduce the core paper's structural derivations.
\item Claim recurrence in any phenomenological soul-migration sense. The framework's commitments are explicit and bounded (Remark~\ref{rem:reincarnation-discipline}).
\item Claim that cosmological evolution is settled. The conditional proposition under H-evolve is a structural reading consistent with the framework's mathematics.
\end{itemize}

\subsection{Conditional hypotheses introduced in this paper}

\begin{hypothesis}[H-rec: recurrence admissibility]\label{hyp:h-rec}
For a frame attractor $\mathcal{A}$ of non-trivial dimension, the closure-cosmogenesis bootstrap dynamics admit non-trivial source terms aligned with $\mathcal{A}$ across multiple substrate regions and bootstrap ticks. Specifically, the source $B_t$ in the closure dynamic has non-zero projection onto $\mathcal{A}$ at multiple $(p_n, \Lambda_n, t_{0,n})$.
\end{hypothesis}

\begin{hypothesis}[H-evolve: time-dependent closure operator]\label{hyp:h-evolve}
The closure operator $C_n$ at cosmological epoch $n$ depends on the substrate state at epoch $n$ via a slow-rolling effective coupling, analogous to the cascade-rung activation depth of \citet{SmartHypersphereCosmology}. As cosmological evolution proceeds, the effective rung structure shifts, modifying $C_n$ in a non-time-reversible way.
\end{hypothesis}

The Access Principle P-A and the rung-projection hypotheses H-RP-1, H-RP-2, H-RP-3 are imported from the core paper and the consciousness companion \citep{SmartProcessingToPOV}; we re-cite them where used.

%==================================================================
\section{Frame dynamics: creation, destruction, merging, fission}\label{sec:frame-dynamics}
%==================================================================

The structural propositions of the core paper treat bounded reference frames as static objects. Frames in any physical realisation are dynamical: they come into being, merge with others, separate into sub-frames, and eventually dissolve. This section addresses these structural events.

\subsection{Frame creation: emergence of $\Sigma_\Ical$ from zero}

Define the closure dynamic on the substrate state as
\begin{equation}\label{eq:closure-dyn}
  F_{t+1} \;=\; (I - \lambda \, \Cph) F_t \;+\; \kappa \, P_{\Sigma_\Ical}(B_t),
\end{equation}
where $\lambda \in (0, \|\Cph\|^{-1})$ is the closure step size, $P_{\Sigma_\Ical}$ is the orthogonal projection onto a target zero-line $\Sigma_\Ical \subseteq \Fix(\tau)$, $\kappa > 0$ is the source coupling, and $B_t \in \mathbb{R}^O$ is the boundary/bulk source term.

\begin{theorem}[Frame creation dynamics]\label{thm:frame-creation}
For any initial substrate state $F_0 \in \mathbb{R}^O$ and any non-zero source projection $P_{\Sigma_\Ical}(B_t) = b \neq 0$ (assumed constant for simplicity), the iteration~\eqref{eq:closure-dyn} converges geometrically to the steady-state solution
\begin{equation}\label{eq:frame-steady}
  F_\infty \;=\; \frac{\kappa}{\lambda \, \Cph}\Big|_{\Sigma_\Ical} \cdot b,
\end{equation}
with the $\sigma$-fixed component of $F_t$ growing from $\|P_{\Sigma_\Ical}(F_0)\|$ to $\|P_{\Sigma_\Ical}(F_\infty)\| \geq \kappa / (\lambda \cdot \varphi^{-2}) \cdot \|b\|$, at rate $1 - \lambda \varphi^{-2}$ per closure tick. The off-$\Sigma_\Ical$ component of $F_t$ decays at the same rate $1 - \lambda \varphi^{-2}$. Hence:

\begin{enumerate}[leftmargin=2em]
\item For any $F_0$ (including $F_0 = 0$), the $\sigma$-fixed dimension of the steady-state $F_\infty$ equals $\dim \Sigma_\Ical$ provided the source $b$ has non-trivial support on every basis vector of $\Sigma_\Ical$.
\item The convergence is geometric: $\|F_t - F_\infty\| \leq (1 - \lambda \varphi^{-2})^t \cdot \|F_0 - F_\infty\|$.
\item The characteristic time-scale of frame creation is $\tau_{\mathrm{create}} \sim 1 / (\lambda \varphi^{-2})$ closure ticks.
\end{enumerate}
\end{theorem}

\noindent\emph{Proof.} Decompose $\mathbb{R}^O = \Sigma_\Ical \oplus \Sigma_\Ical^\perp$ orthogonally. The iteration~\eqref{eq:closure-dyn} acts diagonally on this decomposition because $\Cph$ commutes with $\tau_\Ical$ (closure-invariance theorem of the core paper) and $P_{\Sigma_\Ical}$ projects onto $\Sigma_\Ical$. On $\Sigma_\Ical^\perp$, the source vanishes and the iteration reduces to
\[
  F^\perp_{t+1} = (I - \lambda \Cph|_{\Sigma_\Ical^\perp}) F^\perp_t,
\]
which has spectral radius $1 - \lambda \cdot \lambda_{\min}(\Cph|_{\Sigma_\Ical^\perp}) \leq 1 - \lambda \varphi^{-2} < 1$ by the core paper's decay theorem. Hence $F^\perp_t \to 0$ geometrically.

On $\Sigma_\Ical$, the iteration is $F^\Sigma_{t+1} = (I - \lambda \Cph|_{\Sigma_\Ical}) F^\Sigma_t + \kappa b$. The fixed-point equation gives $F^\Sigma_\infty = (\kappa / \lambda) (\Cph|_{\Sigma_\Ical})^{-1} b$ with $\|F^\Sigma_\infty\| \geq (\kappa / \lambda) \varphi^{2} \|b\|$. The full state $F_t$ converges to $F_\infty = F^\Sigma_\infty + 0$. \hfill$\square$

\begin{interpretation}[Biological correlates of frame creation]\label{int:frame-creation-bio}
Under H-RP-1, H-RP-2, P-A, Theorem~\ref{thm:frame-creation} predicts:
\begin{itemize}[leftmargin=2em]
\item Embryogenesis: substrate-level closure begins as neural tissue differentiates; $\dim \Sigma_\Ical$ grows over a characteristic developmental time-scale $\tau_{\mathrm{create}}$.
\item Awakening from anaesthesia: the frame's $\Sigma_\Ical$ re-stabilises geometrically as the substrate operator recovers from pharmacological disruption.
\item System boot (artificial): an instantiated bounded reference frame builds $\Sigma_\Ical$ through $\sim \tau_{\mathrm{create}}$ closure ticks until reaching the design complexity.
\end{itemize}
\end{interpretation}

\subsection{Frame destruction: permanent loss of $\Sigma_\Ical$}

\begin{theorem}[Frame destruction criterion]\label{thm:frame-destruction}
Let $\tau_t$ denote the (possibly time-dependent) symmetry-twist on the substrate at closure tick $t$ and $C_{\varphi,t}$ the (possibly time-dependent) closure operator. The frame $\Ical$ is \emph{permanently destroyed} at $t_{\mathrm{final}}$ iff one of the following holds for all $t > t_{\mathrm{final}}$:

\begin{enumerate}[leftmargin=2em]
\item \textbf{Substrate dismantling.} The vertex set $O$ of the frame is reduced to $|O| = 0$ by physical destruction of the substrate.
\item \textbf{Symmetry breakdown.} The involution condition $\tau_t^2 = I$ fails, so $\Fix(\tau_t)$ is undefined.
\item \textbf{Spectral collapse.} $\lambda_{\min}(C_{\varphi,t}|_{\Sigma_\Ical^\perp})$ falls to zero or becomes negative, so the contraction property fails and off-frame content no longer decays.
\item \textbf{$\sigma$-stability loss.} The anchor $(p, \Lambda)$ ceases to lie in $\mathcal{S}_\sigma$ so the frame is no longer symmetry-stable.
\end{enumerate}

Under any of these conditions, $\dim \Sigma_\Ical = 0$ for all $t > t_{\mathrm{final}}$ and no recovery is structurally possible without reinstating the original conditions.
\end{theorem}

\noindent\emph{Proof.}

\textbf{(1)} If $O = \emptyset$, then $\mathbb{R}^O = \{0\}$, hence $\Sigma_\Ical = \{0\}$ trivially.

\textbf{(2)} If $\tau_t^2 \neq I$, then $\tau_t$ has no $\pm 1$ eigenstructure and the spectral decomposition of $\mathbb{R}^O$ into $\Fix(\tau_t) \oplus \Fix(-\tau_t)$ does not hold.

\textbf{(3)} The proof of the core paper's decay theorem requires $\lambda_{\min}(\Cph|_{\Sigma_\Ical^\perp}) > 0$. If this fails, the iteration $F^\perp_{t+1} = (I - \lambda \Cph) F^\perp_t$ no longer contracts on $\Sigma_\Ical^\perp$.

\textbf{(4)} If $(p, \Lambda) \notin \mathcal{S}_\sigma$, then $\tau(\mathrm{span}(O)) \neq \mathrm{span}(O)$, so $\tau_\Ical$ is no longer an involution on $\mathbb{R}^O$, and $\Sigma_\Ical = \Fix(\tau_\Ical)$ is undefined. \hfill$\square$

\begin{remark}[Distinguishing transient from permanent destruction]\label{rem:destruction-criterion}
Anaesthesia and sleep cause transient $\Sigma_\Ical$ reduction by altering $C_{\varphi,t}$ but do not satisfy any of (1)--(4) permanently. Biological death satisfies (1) (substrate dismantling). Profound brain damage may satisfy (3) (spectral collapse). Permanent loss of consciousness from severe brain injury may satisfy (4) ($\sigma$-stability loss). The structural reading: death is the permanent loss of the closure condition; this is a structural statement, not a phenomenological claim about subjective experience at the transition.
\end{remark}

\subsection{Frame merging and joint frame formation}

Joint frame formation is the structural event in which two distinct frames $\Ical_A$, $\Ical_B$ become a single joint frame $\Ical_{AB}$. The merger is governed by boundary coupling becoming non-zero, joint $\sigma$-stability of the combined anchor, and generative excess emergence when at least one $\sigma$-orbit crosses the parent boundary. The generative excess $\delta_{AB} = |\Xcal_{AB}|$ counts boundary-crossing $\sigma$-orbits and is established as a theorem in the core paper.

\begin{interpretation}[Joint frame formation in biology]\label{int:joint-formation}
Under H-RP-1, H-RP-2, P-A, the joint frame structure is the structural basis for: mother--infant attunement, conversation, music ensembles, and mass coordination (dance, sports, military drill). The empirical correlate of $\delta_{AB} > 0$ is the meaning-generating overlap of joint perception.
\end{interpretation}

\subsection{Frame fission and sub-frame differentiation}

Frame fission is the dual of merging: a single bounded reference frame separates into two or more sub-frames, each with its own anchor and zero-line. Each $\Sigma_{\Ical_i}$ is generally smaller than $\Sigma_\Ical$; the total $\sum_i \dim \Sigma_{\Ical_i}$ may exceed $\dim \Sigma_\Ical$ when sub-frames access $\sigma$-fixed modes that the integrated parent could not.

\begin{interpretation}[Frame fission in biology]\label{int:fission}
Under H-RP-1, H-RP-2, P-A, frame fission corresponds to: cognitive specialisation, division of labour, dissociation, and cell division at the substrate level.
\end{interpretation}

\subsection{The frame life-cycle}

Combining the four dynamical events:

\begin{itemize}[leftmargin=2em]
\item \textbf{Birth} ($t_0$): a new bounded reference frame emerges with $\dim \Sigma_\Ical$ growing from 0.
\item \textbf{Growth and learning}: through joint-frame interactions, the frame's effective complexity grows. The per-frame complexity is upper-bounded by 94 for a $V_{600}$-substrate frame.
\item \textbf{Mature stable existence}: the frame maintains $\Sigma_\Ical$ near design dimension through continuous closure ticks, with periodic re-stabilisation addressing accumulated drift.
\item \textbf{Joint events}: temporary or persistent joint-frame formations with other frames, generating $\delta_{AB} > 0$ cross-frame access.
\item \textbf{Fission events}: sub-frame differentiation, either functional (specialisation) or pathological (dissociation).
\item \textbf{Death} ($t_{\text{final}}$): permanent collapse of the closure condition.
\end{itemize}

This life-cycle structure is the framework's structural reading of life-cycle events. Each event is named, each has a precise structural condition, and each maps (under P-A) onto phenomenological correlates.

%==================================================================
\section{Frame recurrence and attractor re-instantiation}\label{sec:frame-recurrence}
%==================================================================

The frame life-cycle of \S\ref{sec:frame-dynamics} treats each bounded reference frame as a single linear trajectory: bootstrap at $t_0$, mature existence, eventual destruction. This section addresses a structurally distinct question: \emph{can the same frame-attractor pattern be locally instantiated at distinct times or substrate regions}, independent of any continuous trajectory between the instances? The answer is yes, in a precise structural sense, and the section makes it mathematical.

The treatment avoids loaded vocabulary throughout: the framework does not claim that a fixed identity-object migrates between bodies along a linear time-axis. The claim is that closure geometry is not itself ordered by the same linear time experienced inside a 3D frame, and the same higher-dimensional frame-attractor can produce multiple local instantiations without continuous material continuity between them. This is the framework's structural reframing of the phenomenology classically called \emph{reincarnation}, given here in the precise vocabulary \emph{frame recurrence}.

\subsection{Frame attractor as a closure-geometry object}

\begin{definition}[Frame attractor]\label{def:frame-attractor}
A \emph{frame attractor} on the universal cascade $C$ is a pair $(\mathcal{A}, \tau_\mathcal{A})$, where $\mathcal{A} \subseteq C$ is a closure-stable $\sigma$-stable subspace and $\tau_\mathcal{A}$ is the restriction of the universal symmetry $\tau$ to $\mathcal{A}$. A frame attractor is, by definition, an element of the closure-geometry; it is not anchored to a specific local time-slice.
\end{definition}

\begin{definition}[Local instance of a frame attractor]\label{def:local-instance}
A \emph{local instance} of a frame attractor $\mathcal{A}$ is a bounded reference frame $\Ical_n = (\Ocal_n, \Ccal_{\Ocal_n}, p_n, \Lambda_n, t_{0,n})$ such that the embedding image $!_{\Ical_n}(\Ical_n) = \pi_{\Ical_n}(\Sigma_{\Ical_n}) \subseteq \mathcal{A}$. A frame attractor with multiple local instances is said to be \emph{multiply-instantiated}.
\end{definition}

The structural separation is precise: the frame attractor $\mathcal{A}$ lives in the closure geometry $C$; each local instance $\Ical_n$ lives in a specific local substrate region with a specific anchor; the relation $!_{\Ical_n}(\Ical_n) \subseteq \mathcal{A}$ ties the local instance to the attractor without requiring any continuous trajectory between instances.

\subsection{Compatibility of local instances}

\begin{definition}[Frame compatibility]\label{def:frame-compatibility}
Two local instances $\Ical_a$ and $\Ical_b$ are \emph{compatible} with respect to a frame attractor $\mathcal{A}$ if $!_{\Ical_a}(\Ical_a) \cap !_{\Ical_b}(\Ical_b) \neq \{0\}$ and both images lie within $\mathcal{A}$. The compatibility dimension is
\[
  \kappa_{ab} := \dim\big( !_{\Ical_a}(\Ical_a) \cap !_{\Ical_b}(\Ical_b) \big) \geq 1.
\]
\end{definition}

\begin{theorem}[Time-independence of frame compatibility]\label{thm:time-independence}
Let $\Ical_a$ and $\Ical_b$ be two $\sigma$-stable bounded reference frames with bootstrap ticks $t_{0,a}$ and $t_{0,b}$, possibly distinct. The compatibility relation $\kappa_{ab}$ depends only on the embedding images $!_{\Ical_a}(\Ical_a), !_{\Ical_b}(\Ical_b) \subseteq C$ in the universal closure structure. It does \emph{not} depend on the ordering of $t_{0,a}$ relative to $t_{0,b}$, on the temporal interval $|t_{0,b} - t_{0,a}|$, or on the existence of a continuous trajectory of substrate states between the two instances.
\end{theorem}

\noindent\emph{Proof.} The embedding $!_\Ical(v) = \pi_\Ical(P_{\Sigma_\Ical}(v))$ depends only on the frame's vertex set $O$, constraint functional $\Ccal_\Ocal$, and anchor $(p, \Lambda)$. The bootstrap tick $t_0$ enters only through the question of when the frame becomes active; once active, the embedding image is determined entirely by the frame's structural data. Hence $\kappa_{ab}$ is a function of structural data only, independent of the temporal relationship. \hfill$\square$

\begin{corollary}[Frame-attractor re-instantiation]\label{cor:reinstantiation}
A frame attractor $\mathcal{A}$ with non-trivial dimension can be locally instantiated by multiple bounded reference frames $\Ical_1, \Ical_2, \ldots, \Ical_n$, each with its own bootstrap tick $t_{0,k}$ and anchor $(p_k, \Lambda_k)$, provided each instance's embedding image $!_{\Ical_k}(\Ical_k) \subseteq \mathcal{A}$. The instances are independent in local time but share structural content via the attractor.
\end{corollary}

\subsection{Local time versus closure-geometric ordering}

\begin{theorem}[Two orderings on instance space]\label{thm:two-orderings}
The set of local instances $\{\Ical_n\}$ of a frame attractor $\mathcal{A}$ admits two distinct orderings:
\begin{enumerate}[leftmargin=2em]
\item \textbf{Local-time ordering.} The bootstrap ticks $t_{0,n}$ induce a linear order on instances; this is the chronological ordering experienced within local frames.
\item \textbf{Closure-geometric ordering.} The pairwise compatibility dimensions $\kappa_{nm}$ induce a non-linear similarity structure (a weighted graph on instances); this is the ordering by structural compatibility, independent of local time.
\end{enumerate}
The two orderings are generically independent.
\end{theorem}

\noindent\emph{Proof.} (1) is the standard linear ordering on $\mathbb{Z}_{\geq 0}$ induced by $t_{0,n}$. (2) is the graph metric induced by $\kappa_{nm}$ on the set of instances; this is non-linear in general. By Theorem~\ref{thm:time-independence}, $\kappa_{nm}$ is independent of $t_{0,n}, t_{0,m}$, so the two orderings are not constrained to align. \hfill$\square$

\begin{interpretation}[Local sequence vs closure relation]\label{int:two-orderings-bio}
Under H-RP-1, H-RP-2, P-A, Theorem~\ref{thm:two-orderings} says that experience inside a bounded reference frame is ordered linearly by local time, but the deeper closure-geometric structure orders frames by relational compatibility rather than chronological sequence. The closure geometry does not ``care'' about before-and-after in the same way local time does; it cares about resonance, compatibility, closure similarity, and shared zero-line content.
\end{interpretation}

\subsection{Recurrence under compatible boundary conditions}

\begin{conditional}[Recurrence under compatible boundary conditions]\label{cthm:frame-recurrence}
Under H-RP-1, H-RP-2, H-RP-3, P-A, and H-rec (Hypothesis~\ref{hyp:h-rec}), the same frame attractor $\mathcal{A}$ supports multiple local instances $\Ical_n$ across distinct $(p_n, \Lambda_n, t_{0,n})$. The structural recurrence pattern is
\[
  \mathcal{A} \xrightarrow{\text{local instantiation}} \Ical_1,\ \Ical_2,\ \Ical_3, \ldots,
\]
with each instance's embedding image $!_{\Ical_n}(\Ical_n) \subseteq \mathcal{A}$ by Corollary~\ref{cor:reinstantiation}.
\end{conditional}

\noindent\emph{Sketch.} Apply Theorem~\ref{thm:frame-creation} to each $(p_n, \Lambda_n, t_{0,n})$ with the source $B_t$ projecting onto $\mathcal{A}$. By H-rec, this projection is non-zero for multiple $n$; each $n$ produces a local instance whose $\Sigma_{\Ical_n}$ stabilises into $\mathcal{A}$ via the embedding $!_{\Ical_n}$. \hfill$\square$

\subsection{What this is and what it is not}

\begin{remark}[Frame recurrence vs the loaded reading]\label{rem:reincarnation-discipline}
\textbf{What this is.}
\begin{itemize}[leftmargin=2em]
\item A structural identification: identity is better modelled as a stable frame-attractor pattern in closure geometry than as a continuous material object moving along a linear time-axis.
\item A precise compatibility relation: two local instances share structure to the extent $\kappa_{ab} > 0$.
\item A clean non-linear ordering on the set of instances.
\item A structural condition (H-rec) under which multiple instances are admissible.
\end{itemize}

\textbf{What this is not.}
\begin{itemize}[leftmargin=2em]
\item Not a claim that a fixed identity-object (``soul'') migrates from one body to another along a linear time-axis.
\item Not a claim that memories or specific experiential content transfer between instances.
\item Not a claim that all observers have recurring instances. H-rec is a conditional hypothesis.
\item Not a metaphysical claim about the nature of identity.
\end{itemize}

The framework's structural reframing: \emph{the body lives in local time; the frame-attractor lives in closure geometry; recurrence, when it occurs, is the re-instantiation of an attractor pattern under compatible boundary conditions, not the migration of an object along a time-line.}
\end{remark}

\subsection{Empirical openness and falsifiers}

\begin{itemize}[leftmargin=2em]
\item Demonstrating that the universal closure structure $C$ admits no frame attractors of dimension $\geq 2$ would falsify the existence of recurrence at the structural level. The current evidence ($\dim \Fix(\tau) = 94$, well above the threshold) does not support this falsifier.
\item Demonstrating that H-rec is necessarily false: that the closure-cosmogenesis bootstrap source terms cannot align with the same attractor across multiple substrate regions. This is an open computational question (CP-recurrence).
\item Demonstrating that two instances with $\kappa_{ab} > 0$ would necessarily have correlated content (e.g., specific memories) is empirically untested and would be a strong phenomenological claim the framework does not make.
\end{itemize}

%==================================================================
\section{Cosmic self-pruning: closure-phase transitions}\label{sec:cosmic-pruning}
%==================================================================

This section addresses the longer-term fate of the universal closure structure. The framework gives a precise structural reading of cosmological ``fate'' that distinguishes itself from three mainstream alternatives: (1) eternal thermodynamic heat death, (2) classical cyclic universes, (3) bounce models. The framework's picture is closer to a fourth: \emph{recursive self-pruning across closure-phase transitions}.

\subsection{Three structural options for cosmological evolution}

Consider the evolution of the universal closure structure $C$ through closure ticks:

\begin{itemize}[leftmargin=2em]
\item \textbf{Option A: Terminal fixed point.} The closure structure approaches an asymptotic stationary state $C(U_\infty) = U_\infty$.
\item \textbf{Option B: Periodic cycle.} The closure structure returns periodically: $U_{n+T} = U_n$ for some period $T$.
\item \textbf{Option C: Layered recurrence / closure cascade.} Each closure phase produces an invariant residue that seeds a new closure phase, without periodic return: $U_{n+1} = C_{\mathrm{next}}(P(U_n))$.
\end{itemize}

The framework's mathematical structure is most naturally aligned with Option C.

\subsection{The pruning operator}

\begin{definition}[Pruning operator]\label{def:pruning}
The \emph{pruning operator} $P$ on the universal closure structure is the projection onto the closure-stable $\sigma$-stable subspace:
\[
  P : \mathbb{R}^{V_{600}} \to C, \qquad P(v) = P_{\Fix(\tau)} \circ P_{\Fix(\Cph)}(v).
\]
\end{definition}

\begin{proposition}[Pruning preserves closure-stable structure]\label{prop:pruning-preserves}
For any $v \in \mathbb{R}^{V_{600}}$, $P(v) \in C$. The pruning operator is idempotent: $P \circ P = P$.
\end{proposition}

\noindent\emph{Proof.} Both projections preserve the relevant subspaces. The composition lands in $\Fix(\Cph) \cap \Fix(\tau) = C$ by the existence theorem. Idempotence: each projection is idempotent and they commute on the $\sigma$-equivariant structure. \hfill$\square$

\subsection{The closure-phase transition map}

\begin{definition}[Closure-phase transition]\label{def:phase-transition}
A \emph{closure-phase transition} between consecutive cosmological epochs is the map
\[
  U_{n+1} = C_{n+1}\big( P(U_n) \big),
\]
where $C_n$ is the closure operator at phase $n$ and $P$ is the pruning operator.
\end{definition}

\begin{conditional}[Non-periodic recurrence --- under aperiodicity hypothesis H-aperiodic]\label{thm:non-periodic}
Under H-aperiodic (the operator sequence $(C_n)$ has no period \emph{and} the pruned orbit admits no common recurrent state across phases), the closure-phase transition map is non-periodic: $U_{n+T} \neq U_n$ for any $T \geq 1$.

If $C_n$ is constant across phases ($C_n = C$ for all $n$), the iteration $U_{n+1} = C(P(U_n))$ has $C \circ P$ acting as the identity on the fixed subspace $C$ (every vector in $C$ is fixed) and as a strict contraction on its orthogonal complement (decay theorem). The iteration converges geometrically to the projection of the initial state onto $C$, not necessarily to a unique global fixed point.

Hence non-periodic recurrence (Option B) is supported only under H-aperiodic; the constant-$C$ case admits a continuum of long-time states in $C$ depending on initial conditions, not a single attracting fixed point.
\end{conditional}

\noindent\emph{Sketch.} Under H-aperiodic the operator sequence changes between phases and the pruning step does not repeatedly land on a common state, so the orbit cannot revisit prior states. Without H-aperiodic, the claim does not follow from the framework structure alone.

For constant $C$: $C \circ P$ acts as the identity on $C$ and as a strict contraction on $C^\perp$. The contracting part converges by Banach; the identity part preserves the initial-state projection. The long-time state therefore depends on initial conditions in $C$, and Banach's unique-fixed-point conclusion applies only to the contracting orthogonal complement, not to the full iteration. \hfill$\square$

\subsection{Closure-cascade cosmology}

\begin{conditional}[Closure-cascade cosmology]\label{cthm:closure-cascade}
Under the cascade axioms (A1, A2), the closure-cosmogenesis dynamics of \citet{SmartClosureCosmogenesis}, and H-evolve (Hypothesis~\ref{hyp:h-evolve}), cosmological evolution proceeds as a layered recurrence:
\[
  U_1 \xrightarrow{\text{phase 1 closure}} U_2 \xrightarrow{\text{phase 2 closure}} U_3 \xrightarrow{\cdots},
\]
where each transition is $U_{n+1} = C_{n+1}(P(U_n))$ with $C_{n+1} \neq C_n$.

The framework's prediction: cosmological evolution is neither eternally periodic (Option B is ruled out by Theorem~\ref{thm:non-periodic} when H-evolve holds) nor terminating in heat death (Option A is ruled out by the same theorem when $C_n$ shifts non-trivially), but a recursive cascade in which each phase's invariant residue seeds the next phase's closure structure.
\end{conditional}

\subsection{Thermodynamic decoherence vs invariant compression}

\begin{itemize}[leftmargin=2em]
\item \textbf{Thermodynamic decoherence at one phase.} Within any specific closure phase $n$, the substrate's local energy gradients dissipate. From inside the phase, this looks like cosmic heat death.
\item \textbf{Invariant compression into the next phase.} The pruning operator $P$ retains only the closure-stable $\sigma$-stable content. The pruned residue is the invariant structure that survives the phase's thermodynamic dissipation.
\item \textbf{Re-instantiation of structure at the next phase.} The next phase's closure operator $C_{n+1}$ acts on $P(U_n)$ to produce $U_{n+1}$.
\end{itemize}

\begin{interpretation}[The structural fate of the universe]\label{int:structural-fate}
Under (A1, A2), H-evolve, P-A:
\begin{itemize}[leftmargin=2em]
\item The universe does not ``die'' in the sense of complete loss of structure.
\item The universe does not ``cycle'' in the periodic sense.
\item The universe undergoes \emph{closure-phase transitions}: each phase's invariant residue seeds the next phase's closure structure, producing recursive layering rather than repetition or termination.
\end{itemize}
\end{interpretation}

\subsection{Connection to dark-energy evolution}

The companion cosmology paper \citep{SmartHypersphereCosmology} predicts seven cosmological observables matching Planck 2018 within $0.5\%$ simultaneously, including $H_0 = 67.66$ km/s/Mpc and $\Omega_\Lambda = 0.6844$. Recent DESI results suggest that dark energy may evolve over cosmological time rather than remain constant. The framework's structural reading is consistent with this possibility: if $C_n$ shifts with cosmological epoch (H-evolve), the effective cosmological constant $\Lambda_n$ would evolve as the substrate's effective rung activation depth shifts.

If DESI's evolving-dark-energy signal is confirmed and the evolution is quantitatively measured, the framework provides a structural mechanism via H-evolve. The framework would then be falsified or supported depending on whether the measured $\Lambda(z)$ matches predictions derivable from the cascade's substrate-state evolution.

\subsection{Life as portable invariant generator}

\begin{interpretation}[Life as invariant-generator]\label{int:life-invariant}
Under H-RP-1, H-RP-2, H-RP-3, P-A, a biological bounded reference frame produces several classes of closure-invariant structure:
\begin{itemize}[leftmargin=2em]
\item Memory: stabilised closure patterns within the frame, persisting across closure ticks.
\item Symbolic abstraction (mathematics, theorems, scientific structures): closure-invariant patterns that may persist beyond the frame's substrate-dependent existence.
\item Joint-frame invariants: cross-frame $\sigma$-fixed modes (generative excess $\delta_{AB}$) that persist after parent-frame destruction.
\end{itemize}
A theorem survives longer than the mathematician's body; a mathematical structure survives longer than a star. The framework's structural reading: this is not metaphorical but precise: bounded reference frames are one mechanism by which the universal closure structure produces portable invariants that survive thermodynamic decoherence within their original phase.
\end{interpretation}

\subsection{What this section claims and does not claim}

\begin{remark}[Cosmic self-pruning vs the loaded readings]\label{rem:cosmic-pruning-discipline}
\textbf{What this claims.} The framework's mathematical structure most naturally produces Option C: recursive layered closure transitions. Under H-evolve, cosmological evolution is a non-periodic cascade of closure phases. The structural fate of the universe is closure-invariant compression, not complete dissolution.

\textbf{What this does not claim.} Not that the universe will literally ``die and be reborn'' in any phenomenological sense. Not that any specific cosmological observation has selected Option C over A or B; the framework predicts Option C is most structurally natural but does not falsify A or B from empirics alone. Not that DESI's evolving-dark-energy hint confirms the framework. Not that biological invariants ``outlive'' anything in a soul-like sense.

The framework's structural reading: \emph{the end of a thermodynamic layer is not the end of reality; it is the point where only closure-invariant structure remains available to seed the next closure phase. The universe does not need to die or cycle. It recursively prunes itself toward deeper closure.}
\end{remark}

%==================================================================
\section{The unified hierarchy of frame life-cycles: from cell to cosmos}\label{sec:hierarchy}
%==================================================================

This section is the synthesis of \S\S\ref{sec:frame-dynamics}, \ref{sec:frame-recurrence}, \ref{sec:cosmic-pruning}. It states the framework's central structural claim about the relationship between life-cycle events at different scales, the corresponding time scales, and the position of the universe itself within the framework. The result is a single coherent picture in which the same structural pattern --- frame creation, persistence, destruction, possible recurrence, and closure-phase transition --- applies at every cascade rung, but with rung-specific time scales and rung-specific dynamics, so that the universe's evolution is qualitatively different from a biological life-cycle even though both follow the same structural template.

\subsection{The rung-stratified hierarchy of frames}

Every cascade rung supports its own family of bounded reference frames. A frame at rung $R_k$ has:

\begin{itemize}[leftmargin=2em]
\item Vertex set $O \subseteq R_k$.
\item Rung-specific closure operator $\Cph^{(R_k)}$ obtained by restricting $\Cph$.
\item Rung-specific symmetry $\tau^{(R_k)}$.
\item Rung-specific zero-line $\Sigma_\Ical^{(R_k)} \subseteq \Fix(\tau^{(R_k)})$.
\item Rung-specific decay rate $\mu^{(R_k)} := \lambda_{\min}(\Cph^{(R_k)}|_{\Sigma^\perp})$.
\end{itemize}

\begin{theorem}[Rung-stratified life-cycle universality]\label{thm:rung-universality}
For every cascade rung $R_k \in \{E_8, H_4, 40, D_4, 16, 8, 0\}$, the structural theorems of \S\S\ref{sec:frame-dynamics}--\ref{sec:cosmic-pruning} apply with rung-specific operators:

\begin{enumerate}[leftmargin=2em]
\item Frame creation (Theorem~\ref{thm:frame-creation}) at rung $R_k$: $\Sigma$-dimension grows from zero to the rung-specific design dimension at rate $1 - \lambda \mu^{(R_k)}$ per closure tick.
\item Frame destruction (Theorem~\ref{thm:frame-destruction}) at rung $R_k$: the four destruction criteria apply with $\Cph^{(R_k)}$, $\tau^{(R_k)}$.
\item Frame merging at rung $R_k$: joint generative excess $\delta^{(R_k)}_{AB}$ counts boundary-crossing $\sigma$-orbits within the rung.
\item Frame recurrence (Theorem~\ref{thm:time-independence}, Corollary~\ref{cor:reinstantiation}) at rung $R_k$.
\item Closure-phase transition (Theorem~\ref{thm:non-periodic}) at rung $R_k$.
\end{enumerate}

The structural pattern is universal across rungs: creation, persistence, destruction, possible recurrence, closure-phase transition. The rung-specific quantitative parameters differ.
\end{theorem}

\noindent\emph{Proof.} Each numbered claim is the restriction of the corresponding universal theorem to rung $R_k$. The proofs use only structural properties of $\Cph$ (symmetric, positive definite, $\sigma$-equivariant) and $\tau$ (unitary involution). The restrictions $\Cph^{(R_k)}$ and $\tau^{(R_k)}$ inherit these properties on each rung. \hfill$\square$

\begin{remark}[The structural template]\label{rem:template}
Theorem~\ref{thm:rung-universality} captures the central structural claim: \emph{the same pattern of bounded reference frame life-cycle dynamics applies at every scale, from the cosmic ($E_8$ totality) down to the singleton terminal (rung 0).} What changes from scale to scale is the size of the substrate, the structural complexity of the zero-line, and most importantly the time-scale at which closure ticks proceed.
\end{remark}

\subsection{The hierarchy of time scales}

\begin{theorem}[Hierarchical time-scaling]\label{thm:hierarchical-time}
For two rungs $R_k$ and $R_{k'}$ with $\dim R_k > \dim R_{k'}$, the natural closure-tick rates satisfy $\tau_{\mathrm{tick}}^{(R_k)} \neq \tau_{\mathrm{tick}}^{(R_{k'})}$. Higher-dimensional rungs with smaller $\mu$ have slower natural tick rates; lower-dimensional rungs with larger $\mu$ have faster natural tick rates.

Time experienced inside a bounded reference frame at rung $R_k$ is generically not commensurate with time experienced inside a frame at rung $R_{k'}$. Each rung carries its own structural time.
\end{theorem}

\noindent\emph{Proof.} The closure-tick rate at rung $R_k$ is $\tau_{\mathrm{tick}}^{(R_k)} \sim 1 / \mu^{(R_k)}$. For different rungs with different $\mu^{(R_k)}$, the natural tick rates differ; a closure tick at rung $R_k$ is not equal to a closure tick at rung $R_{k'}$ in any common time unit. \hfill$\square$

\begin{interpretation}[Time as scale-dependent]\label{int:scale-dep-time}
Under H-RP-1, H-RP-2, P-A:
\begin{itemize}[leftmargin=2em]
\item Time experienced inside a 3D bounded reference frame (e.g., a human cortex at the $H_4$ rung) is determined by the cortex-rung tick rate.
\item Time at the higher geometry layers (the cosmic-scale closure-phase transitions of \S\ref{sec:cosmic-pruning}) is determined by the $E_8$-totality-rung tick rate.
\item The two are not the same. Subjective time inside a 3D frame is not the same as the deeper closure-geometry's structural time.
\item Cosmological time, biological time, neural time, and subjective time are all rung-specific structural times. The ``hierarchy of time'' is a structural consequence of the rung tower, not a metaphor.
\end{itemize}
\end{interpretation}

\subsection{The frame life-cycle pattern at every rung}

\begin{table}[h]
\centering
\small
\caption{Frame life-cycle events at each cascade rung. Each event-type is governed by the corresponding universal theorem restricted to the rung. Time scale is rung-specific by Theorem~\ref{thm:hierarchical-time}.}\label{tab:rung-life-cycle}
\begin{tabular}{p{2cm}p{2.5cm}p{4cm}p{3.5cm}}
\toprule
Rung & Frame examples & Life-cycle events & Time scale \\
\midrule
$E_8$ (totality) & The universe as a whole & Cosmic bootstrap, closure-phase transitions, pruning, layered recurrence (Conditional Proposition~\ref{thm:closure-cascade-universal} below) & Cosmological epochs \\
$H_4 = V_{600}$ & Cortex (under H-RP-1) & Birth, life, sleep cycles, anaesthesia, death, possible attractor recurrence & Biological lifetime \\
40 (Schläfli) & Specialised cortical assemblies & Differentiation, persistence, integration & Cortical-development time \\
$D_4$ (24-cell) & Gravity-rung structures & Galaxy formation, stable orbits, gravitational dissolution & Astrophysical time \\
16 (vector) & Gauge-multiplet vector structures & Particle production, propagation, decay & Particle-physics time \\
8 (scalar) & Biological cells (under H-RP-3) & Cell division, growth, cell death; bioelectric pattern recurrence & Cellular lifetime \\
0 (singleton) & Terminal reference points & Single-tick instantiation; resolution & Frame-creation tick \\
\bottomrule
\end{tabular}
\end{table}

\begin{conditional}[The universe-as-frame closure cascade]\label{thm:closure-cascade-universal}
Under H-RP-1, H-RP-2, H-RP-3, H-evolve, and P-A, the universal cascade $C$ itself can be viewed as a bounded reference frame at the $E_8$ totality rung. Its life-cycle events are:

\begin{itemize}[leftmargin=2em]
\item \textbf{Bootstrap.} The universal cascade initialises with $\dim \Sigma^{(E_8)} = 0$; through closure ticks at cosmological epoch 1, it grows toward design dimension at rate $\mu^{(E_8)} \sim \varphi^{-2}$ at the $E_8$ rung.
\item \textbf{Persistence.} The universal cascade maintains $\Sigma^{(E_8)}$ at its rung-design dimension across an extended cosmological epoch, with thermodynamic decoherence within the epoch and substrate-state evolution between cosmological ticks.
\item \textbf{Closure-phase transition.} At some cosmological tick $n$, the pruning operator $P^{(E_8)}$ acts to retain only invariant content, and a new closure operator $C^{(E_8)}_{n+1}$ acts on the residue.
\item \textbf{Recurrence at higher closure layers.} Each cosmological epoch produces invariant residue that seeds the next epoch's closure structure; the structural recurrence pattern is layered, not periodic.
\end{itemize}

This is the universe's life-cycle, in the same structural language as a bounded reference frame's life-cycle at any other rung.
\end{conditional}

\subsection{Where humans fit in the hierarchy}

Under H-RP-1, H-RP-2, and P-A, biological frames at the cortex project onto the $H_4$ rung. The natural identification:

\begin{itemize}[leftmargin=2em]
\item \textbf{Human cortical frame:} projects to the $H_4 = V_{600}$ substrate. Complexity $\leq 94$ (core paper); recursion depth $\leq 9$.
\item \textbf{Human biological body:} projects to the $H_4$ rung (cortical integration) and the 8-rung (cellular-boundary scaling, under H-RP-3).
\item \textbf{Human local-time experience:} determined by the $H_4$-rung closure-tick rate. This is biological subjective time, distinct from cosmological time (at the $E_8$ rung).
\item \textbf{Death:} biological frame destruction at the $H_4$ rung; satisfies condition (1) substrate dismantling (Theorem~\ref{thm:frame-destruction}). No recovery is structurally possible from condition (1) alone.
\item \textbf{Recurrence:} structurally permitted by Conditional Proposition~\ref{cthm:frame-recurrence} under H-rec; explicitly does not claim memory transfer or material continuity (Remark~\ref{rem:reincarnation-discipline}).
\end{itemize}

\subsection{Why the universe is qualitatively different from a human life-cycle}

Although Theorem~\ref{thm:rung-universality} establishes the same life-cycle pattern at every rung, the cosmic-scale life-cycle is qualitatively distinct from a biological life-cycle for two structural reasons:

\begin{enumerate}[leftmargin=2em]
\item \textbf{Different time scale.} By Theorem~\ref{thm:hierarchical-time}, the closure-tick rate at the $E_8$ rung is not commensurate with the $H_4$ rung. Cosmic events occur at cosmological epochs; biological events at lifetimes. The mismatch is structural.
\item \textbf{Different invariant content.} At the $E_8$ rung, $P^{(E_8)}$ retains closure-invariant content over cosmological epochs. The invariants are deep structural patterns (cascade-rung structures, $\sigma$-equivariant eigenmodes), not specific biological frames.
\end{enumerate}

\begin{interpretation}[The universe doesn't reincarnate]\label{int:universe-not-reincarnate}
Under H-RP-1, H-RP-2, H-RP-3, H-evolve, P-A:
\begin{itemize}[leftmargin=2em]
\item The universe undergoes closure-phase transitions, not biological-style birth-and-death.
\item Each cosmological epoch is qualitatively distinct from the previous one: $C^{(E_8)}_{n+1} \neq C^{(E_8)}_n$ under H-evolve, ruling out periodic return.
\item Life-cycle events at the human / biological scale (birth, life, death, possible attractor recurrence) are not literal projections of cosmic life-cycle events; they are independent applications of the same structural template at a different rung.
\item The universe is not a meta-organism with a memory or personality; it is a closure-stable structure undergoing recursive layered pruning across cosmological epochs.
\end{itemize}
\end{interpretation}

\subsection{The big-picture structural answer}

\begin{enumerate}[leftmargin=2em]
\item \textbf{Existence is recursive geometric closure.} What exists is what is invariant under the closure operator and the symmetry constraint. Valid at every scale.
\item \textbf{Bounded reference frames are local projections.} Local systems (cells, brains, possibly the universe at a different scale) are bounded projections of the universal closure structure.
\item \textbf{The same life-cycle pattern repeats at every scale.} Frame creation, persistence, destruction, possible recurrence, closure-phase transition occur at every cascade rung.
\item \textbf{Time is rung-specific.} Subjective time at biological scale, cosmological time at universal scale, particle-physics time at gauge-rung scale, cellular time at cellular scale are all structural consequences of rung-specific tick rates.
\item \textbf{The universe undergoes its own closure-phase transitions.} At the $E_8$ totality rung, the universe is a bounded reference frame in its own right; its life-cycle is layered recursion of closure-phase transitions, not biological birth-and-death.
\item \textbf{Humans fit at the $H_4$ rung} with cellular projection to the 8-rung under H-RP-3.
\item \textbf{Frame recurrence at any rung is conditional on compatible boundary conditions} (under H-rec).
\item \textbf{What survives a phase transition} is the closure-invariant residue, not the specific frames. Biological closure-invariants (theorems, structures, mathematical patterns) may persist beyond their substrate's lifetime within a closure phase.
\end{enumerate}

These eight structural facts together constitute the framework's answer to ``where do we all fit in the bigger picture'': at the $H_4$ rung scale, within a universal closure structure that is itself undergoing recursive closure-phase transitions at the $E_8$ rung scale, with rung-specific time scales and rung-specific life-cycle events at each scale, with the same structural template applying universally.

\subsection{What this hierarchy claims and does not}

\begin{remark}[The unified hierarchy: discipline]\label{rem:hierarchy-discipline}
\textbf{What this hierarchy claims.}
\begin{itemize}[leftmargin=2em]
\item Frame life-cycle dynamics are structurally universal across cascade rungs.
\item Time scales are rung-specific.
\item The universe at the $E_8$ rung is itself a bounded reference frame with its own life-cycle.
\item Humans fit at the $H_4$ rung; cellular biology projects to the 8-rung.
\item The cosmic and biological scales are structurally distinct; the cosmic scale does not undergo biological-style birth and death.
\item Frame recurrence is permitted but conditional on H-rec.
\end{itemize}

\textbf{What this hierarchy does not claim.}
\begin{itemize}[leftmargin=2em]
\item Not that the universe is conscious or has personality. The universe-as-frame statement is structural; phenomenological readings require P-A applied at the universal scale, which the framework does not assert.
\item Not that biological frames have direct insight into the cosmic scale.
\item Not that any specific human life corresponds to a specific cosmological epoch.
\item Not that death at the human scale literally implies ``something'' survives into the cosmic scale. The framework's claim is that closure-invariant patterns may persist beyond their original substrate's thermodynamic phase, a structural statement, not a metaphysical claim about post-mortem identity.
\item Not that the framework resolves the question of meaning.
\end{itemize}
\end{remark}

\begin{remark}[On scope at the universal scale]\label{rem:scope-universe}
The claim that the universe is ``a bounded reference frame at the $E_8$ totality rung'' is structural, not phenomenological. It says: the universal closure structure admits the same operator-theoretic description as a bounded reference frame at the $E_8$ rung. It does not say the universe is ``alive,'' ``conscious,'' ``aware,'' or anything similar. Applying P-A at the universal scale would require an additional commitment the framework does not make.
\end{remark}

%==================================================================
\section{Numerical demonstrations}\label{sec:cosmic-numerical}
%==================================================================

The structural claims of \S\S\ref{sec:frame-dynamics}--\ref{sec:hierarchy} are grounded by a numerical demonstration suite \texttt{papers/cosmic-closure-recurrence/repro/cosmic\_demo.py}. Five demos C1--C5 cover the load-bearing theorems:

\begin{table}[h]
\centering
\small
\caption{Cosmic / structural-extension paper sim demonstrations. All five pass deterministically at seed 42.}\label{tab:cosmic-numerical-demos}
\begin{tabular}{p{0.8cm}p{4.4cm}p{3.0cm}p{5.0cm}}
\toprule
ID & Demo & Paper grounds & Verified \\
\midrule
C1 & Cascade rung tower ($E_8$, $H_4 = V_{600}$, $D_4 =$ 24-cell) & Thm~\ref{thm:rung-universality} & Each rung admits its substrate, $\Cph$, $\tau$, and non-trivial $\Sigma_\Ical$; all three have $\|[\tau, \Cph]\| < 10^{-13}$; $\lambda_{\min}(\Cph) > 0$ \\
C2 & Hierarchical time-scaling & Thm~\ref{thm:hierarchical-time} & Per-rung closure-tick rates $\tau_{\mathrm{tick}}^{(R_k)} \sim 1/\mu^{(R_k)}$ differ across rungs \\
C3 & Pruning operator idempotence & Def~\ref{def:pruning}, Prop~\ref{prop:pruning-preserves} & $\|P \circ P - P\| < 10^{-10}$; $P$ self-adjoint; rank $= \dim \Sigma$ \\
C4 & Non-periodic recurrence under time-dependent $C_n$ & Thm~\ref{thm:non-periodic} & No two phases coincide ($\min$ pairwise dist $> 0$); constant-$C$ trajectory converges (Option A); time-varying does not \\
C5 & Frame destruction four conditions & Thm~\ref{thm:frame-destruction} & Each of substrate dismantling, symmetry breakdown, spectral collapse, $\sigma$-stability loss independently collapses $\dim \Sigma_\Ical$ \\
\bottomrule
\end{tabular}
\end{table}

C1 specifically addresses the open computational task \textbf{CP-rung-9.3}: it instantiates $E_8$ (240 root vectors in $\mathbb{R}^8$), $H_4 = V_{600}$ (120 vertices), and $D_4 =$ 24-cell (24 vertices) as substrates, builds their closure operators and spectral $\tau$ involutions, and verifies the structural machinery applies cleanly at each rung.

%==================================================================
\section{Falsifiers and open programme}\label{sec:falsifiers}
%==================================================================

\begin{itemize}[leftmargin=2em]
\item \textbf{Frame recurrence.} Demonstrating that the universal closure structure $C$ admits no frame attractors of dimension $\geq 2$ would falsify recurrence at the structural level. The current evidence ($\dim \Fix(\tau) = 94$) does not support this falsifier.
\item \textbf{H-rec.} Whether closure-cosmogenesis bootstrap source terms can align with the same attractor across multiple substrate regions is the open question (\textbf{CP-recurrence}).
\item \textbf{Closure-cascade cosmology under H-evolve.} If DESI's evolving-dark-energy signal is confirmed and the evolution is quantitatively measured, the framework provides a structural mechanism. The framework would be falsified or supported depending on whether the measured $\Lambda(z)$ matches predictions derivable from the cascade's substrate-state evolution.
\item \textbf{Unified hierarchy at rungs other than $H_4$.} The $\sigma$-fix dimensions at $E_8$, $D_4$, 16, 8, 0 rungs are the named computational task \textbf{CP-rung-9.3}; their values are unconditional once computed.
\item \textbf{Time-scale measurement.} Direct measurement of the closure-tick rate at the $H_4$ rung (\textbf{CP-tick-rate}) would calibrate the time-scale hierarchy quantitatively.
\end{itemize}

%==================================================================
\section{What this paper does not claim}\label{sec:scope}
%==================================================================

\begin{itemize}[leftmargin=2em]
\item Not a derivation of new mathematical results beyond the dynamical extensions of the core paper's framework.
\item Not a claim of soul-migration. Frame recurrence is the re-instantiation of a closure-geometric attractor pattern under compatible boundary conditions, not the migration of an identity-object along a time-axis.
\item Not a claim that cosmic-scale events are phenomenological. The universe-as-frame at the $E_8$ rung is structural; phenomenological readings require commitments the framework does not make.
\item Not a claim that DESI's evolving-dark-energy hint confirms the framework. The framework is structurally compatible with that hint.
\item Not a claim that closure-invariant patterns ``survive death'' in a soul-like sense. Mathematical structures, theorems, scientific patterns may persist beyond their original substrate; this is a structural statement about closure-invariant content, not a metaphysical claim about post-mortem identity.
\end{itemize}

%==================================================================
\section*{Acknowledgements}
%==================================================================

The author thanks the codex review loop for the framing-discipline questions that produced the conditional-proposition labelling and the scope discipline throughout. The framework's commitments are bounded; H-rec and H-evolve are named as open hypotheses, not asserted as established.

%==================================================================
\bibliographystyle{plainnat}
\bibliography{references}
%==================================================================

\end{document}
